% Author-review draft. Source checks and pending human verification are recorded in
% reports/2026-10-31-manuscript-howard-review/literature-evidence.md.
\section{Related Work}
\label{sec:related-work}

\subsection{Patient-specific manufacturing and engineering evidence}
% C001; E001,E009,E010
Autologous cell-therapy manufacturing connects resource allocation to a patient's treatment pathway. Wang et al.~\cite{Wang2023} compare centralized and point-of-care supply chains using simulation, assessing cost, fulfillment time, service level, and utilization. Their analysis makes decentralization a conditional design choice, dependent on demand, capacity, and operating rules. Li and White~\cite{Li2023} formulate relocatable capacity and supplier-disruption resilience as a Markov decision process. Tseng et al.~\cite{Tseng2024PatientHealth} incorporate changing patient health into manufacturing and scheduling through SimPAC. These precedents motivate joint attention to resources, logistics, and patient condition; they do not establish that a new controller improves clinical outcomes. In the present synthetic problem, patient identity and service windows constrain specimen routing, while reagent transfers and idle-bioreactor relocation remain distinct resource decisions.

% C002; E011,E012,E013
Published engineering evidence helps identify cost boundaries and process constraints. Ran et al.~\cite{Ran2020} separate fixed and variable costs for academic nonprofit CAR-T production, including qualified personnel, equipment, and consumables. Lopes et al.~\cite{Lopes2020} model automation, scale-out, and supply-chain costs; their scenario assumptions are modeled inputs, not observed failure frequencies. Fitzgerald et al.~\cite{Fitzgerald2022} compare manual and automated expansion of adipose-derived stromal cells and analyze labor by production, quality-control, and release roles. That process is not CAR-T, so its labor savings cannot serve as a transferable productivity multiplier. These studies support explicit engineering assumptions and sensitivity questions, not retrospective calibration of experiments completed with different inputs.

% C003; E014
Clinical observations require a separate interpretation. Nastoupil et al.~\cite{Nastoupil2020} retrospectively follow patients undergoing leukapheresis for standard-of-care axicabtagene ciloleucel, including patients who do not reach infusion. Such pathway evidence motivates tracking attrition alongside cost. It neither identifies the simulator's daily patient-loss mechanism nor establishes the causal effect of a scheduling policy. Clinical attrition, modeled manufacturing failure, and synthetic patient loss must therefore remain distinct outcomes.

\subsection{Learning with operational structure and planning}
% C004; E002,E015,E016
Tseng et al.~\cite{Tseng2025} provide a direct reinforcement-learning antecedent for dynamic capacity planning in decentralized PRM supply chains. Inventory research supplies complementary lessons about benchmarking and uncertainty. Gijsbrechts et al.~\cite{Gijsbrechts2022InventoryDRL} show that carefully tuned deep reinforcement learning can match strong methods for stationary lost-sales, dual-sourcing, and multi-echelon problems, while documenting substantial initial tuning effort. Dehaybe et al.~\cite{Dehaybe2024NonstationaryInventory} condition inventory policies on rolling demand forecasts and train across randomly generated environments. Their results concern specified inventory models and forecast processes, rather than general immunity to distribution shift. Together, these studies motivate strong operational baselines and explicit statements of which uncertainties a learned policy has encountered.

% C005; E017,E018,E019
Hybrid methods retain useful operating structure. Stranieri et al.~\cite{Stranieri2024HybridInventory} assign production quantities to reinforcement learning and shipping decisions to multistage stochastic programming in a two-echelon network. Deep Controlled Learning~\cite{Temizoz2025DeepControlledLearning} uses approximate policy iteration with sequential halving and common random numbers to allocate simulation effort among candidate actions. Residual reinforcement learning instead superposes a learned action correction on a conventional controller, as demonstrated in robot assembly by Johannink et al.~\cite{Johannink2019ResidualRL}. These are distinct mechanisms. They motivate separating a planner's contribution from a learned increment, but do not imply that a correction to candidate costs is an action residual or that a heuristic safeguard guarantees improvement.

% C006; E003,E005,E006,E007
Learning targets also distinguish planning-assisted controllers. Temporal-difference learning updates predictions using differences between successive predictions~\cite{Sutton1988}; DDPG learns a deterministic continuous-action actor and critic~\cite{Lillicrap2016}. TD-MPC combines short-horizon optimization in a learned latent model with a temporally learned terminal value~\cite{Hansen2022}. The cross-entropy method is an optimization procedure~\cite{DeBoer2005}, not evidence that a learned component adds value. A correction fitted to simulated full-tail cost differences under a fixed continuation policy differs from these constructions. Its contribution must be evaluated against the same planner without the correction, with candidate generation and evaluation resources held constant.

\subsection{Graph representations and communication}
% C007; E004,E020
Graph convolution provides a structured mechanism for combining node features with adjacency information~\cite{Kipf2017}. In supply chains, Kotecha and del Rio Chanona~\cite{Kotecha2025GNNMARL} combine graph state representations with multi-agent reinforcement learning and adapt the parameters of heuristic inventory policies. Graph-aware operational control therefore has established antecedents. Its relevance here lies in representing interactions among candidate requests, competing queues, and transferable resources, rather than in a claim that combining graph networks with reinforcement learning is itself new. Neither graph representation nor adjacency creates additional observations or changes the physical feasibility of transfers.

% C008; E021
Communication and representation must also be distinguished. Tolstaya et al.~\cite{Tolstaya2020} train decentralized swarm controllers by imitating a centralized expert and exchange information through local, potentially multihop communication. Their study is an imitation-learning precedent for structured communication, not evidence of reinforcement-learning gains in manufacturing. For the present comparison, removing messages tests the contribution of a communication channel. Isolating graph aggregation instead requires a non-graph comparator with the same accessible observations, message content and timing, global summaries, and decision interface. Network-derived planner scores count as information even when no explicit graph message is transmitted.

% C009; E022
More expressive aggregation remains a hypothesis for future work. Brody et al.~\cite{Brody2022GATv2} identify the query-independent ranking limitation of standard graph attention and introduce GATv2's dynamic attention. That expressiveness result motivates testing alternative representations; it does not predict an operational gain, patient-service benefit, or general superiority over a matched non-graph model. Changes to aggregation, information access, and planning should consequently be attributed separately.

\subsection{Decision quality and net service}
% C010; E023,E024
Decision-focused learning evaluates predictions through their downstream consequences. Donti et al.~\cite{Donti2017TaskBased} train probabilistic models against task objectives in stochastic optimization, including an inventory example. Elmachtoub and Grigas~\cite{Elmachtoub2022SmartPredict} define prediction loss through the decision error induced in an optimization problem. These approaches explain why predictive fit alone need not identify a useful controller. They motivate assessing a candidate-cost model through its selected actions, but do not establish that the present training rule implements either method or inherits their theoretical properties.

% C011; E001,E010,E014,E023,E024
The downstream objective also determines what counts as improvement. For patient-specific manufacturing, an operating-cost reduction can accompany fewer completed treatments or more unfinished service. A weighted objective cannot by itself establish a net service benefit: its component outcomes and accounting horizon must remain visible. Accordingly, cost, patient losses, completions, unfinished service, and adverse scenarios answer different questions. A patient-loss penalty is a modeling choice, not a clinically validated monetary valuation, and an average cost advantage cannot substitute for evidence of patient protection.

\subsection{Uncertainty, generalization, and statistical limits}
% C012; E002,E016,E025
Generalization to broader simulation uncertainty is a substantive open question. Packer et al.~\cite{Packer2018Generalization} distinguish in-distribution evaluation from tests outside the training environment distribution. Forecast-conditioned inventory learning and PRM capacity-planning studies likewise motivate specifying the uncertainty process~\cite{Dehaybe2024NonstationaryInventory,Tseng2025}. For this application, future questions include joint changes in arrivals, disruption dependence, processing times, and observation delays. These are prospective tests, not completed validations. Additional random seeds under one simulator law address sampling variability; they do not establish robustness to a different law or clinical transportability.

% C013; E026,E027
Finally, Henderson et al.~\cite{Henderson2018DRLMatters} document variability and reproducibility problems in deep reinforcement learning. Agarwal et al.~\cite{Agarwal2021RLiable} show how few-run point estimates can support unreliable rankings and advocate interval estimates and distribution-aware summaries. In paired simulation comparisons, uncertainty across independently trained models must remain distinguishable from uncertainty across evaluated worlds; repeated worlds do not create additional independent training runs. Small or unstable graph increments, adverse service trade-offs, and inconclusive comparisons therefore constrain the claims of this study rather than justify a universal ranking of architectures.
